\section{Introduction}
\label{sec:introduction}

A local feature attribution can be read as a list of reasons for a classifier's prediction: which features contribute most, in which direction, and by how much. When the attribution is computed by a stochastic approximation, the displayed list is also a numerical estimate. Repeating the explainer while holding the prediction problem fixed can change that estimate. For credit or health applications, understanding this source of variation is a prerequisite to interpreting the display; it is not by itself evidence that the prediction or its explanation is suitable for a consequential decision.

The SHAP framework defines additive attributions and supplies algorithms with different computational properties \citep{lundberg2017unified}. Its original evaluation already relates sampled-estimator error to evaluation count, and subsequent work studies KernelSHAP uncertainty and convergence \citep{covert2021improving}. The practical question here is conditional: for a specified trained tabular model, input, empirical background and computation budget, how often does \emph{one} approximate attribution preserve both the reference vector's magnitude and its salient signed features? This question requires more than checking whether repeated runs look similar. Agreement among approximations, error relative to a common target, and the adequacy of an explanation for a particular use are distinct properties.

We organize the study around three questions. First, how does the frequency of meeting an explicit one-run numerical criterion change with actual prediction work and CPU time? Second, how closely does agreement between independent runs track accuracy against the same-game reference? Third, do reference-defined weak signals, small importance gaps and numerical ties identify conditions in which the assessment changes? The adequacy thresholds are declared study conventions, with component and sensitivity analyses; they are not accepted lending or clinical standards.

The completed study combines three public tabular tasks with logistic regression, random forests and histogram gradient boosting. For each fitted model, a fixed held-out panel and empirical background isolate explainer randomness. Grouped Kernel SHAP and paired permutation sampling are compared with exact enumeration of the same empirical marginal game. Repeated seeds quantify within-input variability, while inference over inputs preserves their clustering. Requested computation settings are accompanied by actual model-input rows and measured time, because background expansion, cached coalitions and exhaustive evaluation can make nominal settings misleading.

The study's contribution is a reproducible assessment of conditional one-run numerical adequacy, with explicit targets, costs, uncertainty and failure conditions. It does not introduce a new Shapley estimator or claim to discover stochastic attribution error. The measured panel shows budget- and criterion-dependent adequacy: lowest-budget Kernel success spans $.7794$ to $1$, whereas all highest-budget conditions meet the broad event in every observed run. Stricter signed-feature recovery and individual-input failures qualify that result. Conclusions remain conditional on the fitted predictors, tasks, backgrounds, implementations and tested computation range.
